%!TEX TS-program = MikTex
%!TEX encoding = UTF-8 Unicode

\documentclass[martgin, line]{article}
\usepackage{helvet}
\usepackage{tabularx}
\usepackage{longtable}

% LAYOUT
%--------------------------------
% Margins
\usepackage{geometry} 
\geometry{a4paper, left=.75in, right=.75in, top=.75in, bottom=.75in}

% Do not indent paragraphs
\setlength\parindent{0in}

% Enable multicolumns
\usepackage{multicol}
\setlength{\columnsep}{-3.5cm}

% MARGIN NOTES
%--------------------------------
\usepackage{marginnote}
\newcommand{\note}[1]{\marginnote{\scriptsize #1}}
\renewcommand*{\raggedleftmarginnote}{}
\setlength{\marginparsep}{7pt}
\reversemarginpar

% HEADINGS
%--------------------------------
\usepackage{sectsty} 
\usepackage[normalem]{ulem} 
%\sectionfont{\rmfamily\mdseries} 
%\subsectionfont{\rmfamily\mdseries\scshape\normalsize} 
%\subsectionfont{\rmfamily\mdseries\normalsize} 
%\subsubsectionfont{\rmfamily\bfseries\upshape\normalsize} 

% PDF SETUP
%--------------------------------
\usepackage{hyperref}
\hypersetup {
  colorlinks, breaklinks, xetex, bookmarks,
  urlcolor=[rgb]{0.1,0.1,0.858},
  citecolor=[rgb]{0.1,0.1,0.858},
  linkcolor=[rgb]{0.1,0.1,0.858}
}

% Make Sans serif
\renewcommand{\familydefault}{\sfdefault}
 
% \hypersetup
% {
%   pdfauthor={Albert Einstein},
%   pdfsubject={Albert Einstein's CV},
%   pdftitle={Albert Einstein's CV},
%   colorlinks, breaklinks, xetex, bookmarks,
%   filecolor=black,
%   urlcolor=[rgb]{0.117,0.682,0.858},
%   linkcolor=[rgb]{0.117,0.682,0.858},
%   linkcolor=[rgb]{0.117,0.682,0.858},
%   citecolor=[rgb]{0.117,0.682,0.858}
% }


\newcommand\doi[1]{
  \href{http://dx.doi.org/#1}
  {doi:#1}
}

\newcommand\journal[1]{\emph{#1}} 

%variable definitions for references
\input{refvars.tex}

% DOCUMENT
%--------------------------------
\begin{document}

\begin{flushright}
  \textit{\today}
\end{flushright}

%\begin{center}
%  \textbf{CURRICULUM VITAE}


\begin{LARGE}
  \textbf{Albert Einstein}
\end{LARGE}
%\end{center}

\section*{Contact}

Institute for Advanced Study\\
1 Einstein Drive\\
Princeton, NJ 08540\\
Phone: (609) 734-8000\\
\href{mailto:einstein@ias.edu}{einstein@ias.edu}\\


\section*{Education}
\noindent
\begin{tabular}{lll}
1905&PhD&
Physics,
University of Zurich\\
1900&Diploma&
Mathematics and Physics,
Swiss Federal Polytechnic (ETH Zurich)\\
\end{tabular}

\section*{Professional Experience}
\noindent
\begin{tabular}{lll}
1933-1955&
Professor of Theoretical Physics&
Institute for Advanced Study, Princeton\\
1917-1933&
Director, Kaiser Wilhelm Institute for Physics&
Kaiser Wilhelm Society, Berlin\\
1914-1917&
Professor of Theoretical Physics&
Humboldt University of Berlin\\
1912-1914&
Professor of Theoretical Physics&
ETH Zurich\\
1909-1911&
Associate Professor of Theoretical Physics&
University of Zurich\\
1902-1909&
Technical Expert, Class III&
Swiss Patent Office, Bern\\
\end{tabular}

\section*{Honors and Awards}
\begin{tabular}{lp{5.5in}}
1921&
Nobel Prize in Physics, Royal Swedish Academy of Sciences
\\
1925&
Copley Medal, Royal Society
\\
1929&
Max Planck Medal, German Physical Society
\\
1921&
Matteucci Medal, Italian Society of Sciences
\\
1920&
Barnard Medal for Meritorious Service to Science, Columbia University
\\
\end{tabular}


\section*{Memberships}

\begin{tabular}{l}
Prussian Academy of Sciences, 1914-1933\\
Royal Society (Foreign Member), 1921-1955\\
National Academy of Sciences (USA), 1942-1955\\
Royal Netherlands Academy of Arts and Sciences, 1920-1955\\
\end{tabular}


\section*{Bibliography and Products of Scholarship}

\begin{footnotesize}$^*$ indicates a UNC or JHSPH student or trainee;
    $^{**}$ indicates an GIDTRP student or trainee;
    $^\dagger$ indicates corresponding authorship;
    $^\ddagger$ indicates equal contribution
  \end{footnotesize}

\subsection*{Peer-Reviewed Publications (\npeer\ total)} 

\begin{enumerate}
  \input{refs.tex}
\end{enumerate}


\subsection*{Patents (1 total)}
\begin{enumerate}
    \item
  Refrigerator with no moving parts (Einstein-Szilard refrigerator)
   (No. US1781541).
    Einstein A,
    Szilard L.
  \textit{granted}
  \end{enumerate}


\subsection*{Published Articles and Editorials not Peer Reviewed (\nnopeer\ total)}

\begin{enumerate}
  \input{refs_norev.tex}
\end{enumerate}


\subsection*{Books \& Chapters (\nchapters\ total)}

\begin{enumerate}
  \input{chapters.tex}
\end{enumerate}


\subsection*{Preprints (\npreprint\ total)}

\begin{enumerate}
  \input{preprints.tex}
\end{enumerate}


\subsection*{Letters (\nletters\ total)}

\begin{enumerate}
  \input{letters.tex}
\end{enumerate}


\subsection*{Posters and Presentations at Scientific Meetings (\npres\ total)}

\begin{enumerate}
  \input{scipres.tex}
\end{enumerate}


\subsection*{Invited Seminars (3 total)}

\setlength{\extrarowheight}{.75em}
\begin{longtable}[l]{lp{3.75in}l}
    1915&
  \parbox[t]{3.75in} { On the Electrodynamics of Moving Bodies\\
    \scriptsize{%
      \textit{Academy Lecture Series}. %
      Prussian Academy of Sciences.%
    }} & Berlin, Germany\\
    1921&
  \parbox[t]{3.75in} { Geometry and Experience\\
    \scriptsize{%
      \textit{Academy Lecture}. %
      Prussian Academy of Sciences.%
    }} & Berlin, Germany\\
    1921&
  \parbox[t]{3.75in} { The Meaning of Relativity\\
    \scriptsize{%
      \textit{Stafford Little Lectures}. %
      Princeton University.%
    }} & Princeton, NJ\\
  \end{longtable}
\setlength{\extrarowheight}{0em}

\subsection*{Software (2 total)}

\begin{enumerate}
    \item
    Einstein A,
    Straus EG  (1944)
  \href{https://www.ias.edu/software/ufes}{Unified Field Equation
Solver}.
  Institute for Advanced Study
    \item
    Einstein A,
    Rosen N  (1936)
  \href{https://www.ias.edu/software/grlens}{Gravitational Lens
Calculator}.
  Institute for Advanced Study
  \end{enumerate}


\subsection*{Dissertation}
\textbf{Einstein A} (1905) A New Determination of Molecular Dimensions.
Department of Physics, University of Zurich


%% End scholarly outputs---------------------------------------------

\section*{Teaching Record}

\subsection*{Classroom Instruction (Past 3 Years)}

\setlength{\extrarowheight}{.75em}
\begin{longtable}[l]{lp{4in}l}
                    \end{longtable}
\setlength{\extrarowheight}{0em}


\subsection*{Advisees}


Ernst Gabor Straus, PhD (advisor)
\hfill 1944-1948\\ 
\textit{Studies in unified field theory}\\
Institute for Advanced Study\\

Nathan Rosen, PhD (advisor)
\hfill 1934-1936\\ 
\textit{Einstein-Rosen bridge and EPR paradox investigations}\\
Institute for Advanced Study\\






\subsection*{Postdoctoral Trainees}

\setlength{\extrarowheight}{.25em}
\begin{tabular}{ll}   
Leopold Infeld&
1936-1938\\
Peter Bergmann&
1936-1941\\
Bruria Kaufman&
1947-1949\\
\end{tabular}
\setlength{\extrarowheight}{0em}

\subsection*{Oral Exam Participation}
\setlength{\extrarowheight}{.25em}
\begin{longtable}[l]{lp{3.75in}l}
  \emph{Final Defense}\\
  \hline
    Nathan Rosen&
  \parbox[t]{3.75in} {
    Physics
    \scriptsize{%
      \\Massachusetts Institute of Technology
    }
  }&
  May 1936\\
    Valentine Bargmann&
  \parbox[t]{3.75in} {
    Physics
    \scriptsize{%
      \\Princeton University
    }
  }&
  January 1943\\
    Ernst Gabor Straus&
  \parbox[t]{3.75in} {
    Mathematics
    \scriptsize{%
      \\Columbia University
    }
  }&
  June 1948\\
    \\
  \emph{School Wide}\\
  \hline
    John Archibald Wheeler&
  \parbox[t]{3.75in} {
    Physics
    \scriptsize{%
      \\Johns Hopkins University
    }
  }&
  April 1938\\
    David Bohm&
  \parbox[t]{3.75in} {
    Physics
    \scriptsize{%
      \\University of California, Berkeley
    }
  }&
  June 1943\\
    \\
  \emph{Departmental}\\
  \hline
    Valentine Bargmann&
  \parbox[t]{3.75in} {
    Physics
    \scriptsize{%
      \\Institute for Advanced Study
    }
  }&
  June 1942\\
    Ernst Gabor Straus&
  \parbox[t]{3.75in} {
    Mathematics
    \scriptsize{%
      \\Institute for Advanced Study
    }
  }&
  March 1948\\
  \end{longtable}
\setlength{\extrarowheight}{0em}

%%=========================END TEACHING


\section*{Contracts and Grants}

\subsection*{Active}


\textbf{Rockefeller Foundation
\hfill 1952-1956\\
Albert Einstein (PI)
\hfill
\$ 50000 (Total Costs)
\\
Generalized Theory of Gravitation}\\
Research program pursuing a complete unified field theory unifying
gravitation and electromagnetism\\
\textbf{Role: Principal Investigator}
 80\% FTE\\
 RF-UF-1952\\




\subsection*{Completed}

\textbf{Institute for Advanced Study
 (institutional)
\hfill 1933-1955\\
Albert Einstein (PI)
\hfill
\$ 25000 (Total Costs)
\\
Unified Field Theory Research}\\
\textbf{Role: PI}
\\
 IAS-UFT-01\\


\textbf{University of Zurich
 (institutional)
\hfill 1905-1908\\
Albert Einstein (PI)
\hfill
\$ 1200 (Total Costs)
\\
Investigation of Brownian Motion and Molecular Theory}\\
\textbf{Role: PI}
\\
 UZH-BM-1905\\






\section*{Professional Service}

\subsection*{Departmental, School and University}

\subsubsection*{Committee Memberships}

\setlength{\extrarowheight}{.75em}
\begin{longtable}[l]{lp{3.5in}l}
    1922-1932& 
  \parbox[t]{3.5in} { International Committee on Intellectual
Cooperation\\
    \scriptsize{%
      League of Nations
    }} &  Member \\
    1925-1955& 
  \parbox[t]{3.5in} { Board of Governors\\
    \scriptsize{%
      Hebrew University of Jerusalem
    }} &  Governor \\
  \end{longtable}
\setlength{\extrarowheight}{0em}

\subsubsection*{Other  Service}

\setlength{\extrarowheight}{.75em}
\begin{longtable}[l]{lp{5.5in}}
    1940-1945& 
  \parbox[t]{5.5in} { 
    Chaired faculty hiring committee for theoretical physics\\
    \scriptsize{%
      School of Mathematics, IAS
    }
  } \\
    1935-1955& 
  \parbox[t]{5.5in} { 
    Organized weekly physics colloquium series\\
    \scriptsize{%
      School of Mathematics, IAS
    }
  } \\
  \end{longtable}
\setlength{\extrarowheight}{0em}

%\subsection*{School}

%\subsection*{University}

\subsection*{National and International}

\subsubsection*{Advisory Panels}
\setlength{\extrarowheight}{.75em}
\begin{tabular}{lp{5.5in}}   
1946-1948&
\parbox[t]{5.5in}{Chair, Emergency Committee of Atomic Scientists\\
  \textit{Federation of American Scientists}}\\
1925-1955&
\parbox[t]{5.5in}{Member, Scientific Advisory Board\\
  \textit{Hebrew University of Jerusalem}}\\
\end{tabular}
\setlength{\extrarowheight}{0em}

\subsubsection*{Grant Review Panels}
\setlength{\extrarowheight}{.75em}
\begin{tabular}{lp{5.5in}}   
1934&
\parbox[t]{5.5in}{
   
  Physics Fellowship Selection Committee\\
  \textit{Rockefeller Foundation}
   -- \textit{Reviewer}
   -- \textit{RF-1934}
  }\\
\end{tabular}
\setlength{\extrarowheight}{0em}


\subsubsection*{Organized Sessions and Round Tables}
\setlength{\extrarowheight}{.75em}
\begin{tabular}{lp{5.5in}}   
1927&
\parbox[t]{5.5in}{
    Fifth Solvay Conference on Electrons and Photons\\
  \textit{Solvay Conference}
   -- Invited Participant 
  }\\
1911&
\parbox[t]{5.5in}{
    First Solvay Conference on Radiation and the Quanta\\
  \textit{Solvay Conference}
   -- Invited Participant 
  }\\
\end{tabular}
\setlength{\extrarowheight}{0em}


\subsubsection*{Chaired Sessions}

\setlength{\extrarowheight}{.25em}
\begin{longtable}[l]{lll}
    1930&
  Quantum Mechanics Discussion&
  Sixth Solvay Conference\\
  \end{longtable}
\setlength{\extrarowheight}{0em}

\subsubsection*{Editorial Activities}

\underline{Editor}

\begin{tabular}{lll}
1906-1955&
Annalen der Physik
\\
\end{tabular}


\underline{Associate Editor}

\begin{tabular}{lll}
1909-1914&
Annalen der Physik (Associate Editor)
\\
1920-1933&
Zeitschrift fur Physik
\\
\end{tabular}

\underline{Other Editorial Activities}

\begin{tabular}{lll}
1935-1955&
Physical Review&
Advisory Board
\\
1936&
Journal of the Franklin Institute&
Guest Editor, Special Issue on Unified Field Theory
\\
\end{tabular}

\subsection*{Peer Review Activities}
Reviewed for:
Physical Review;
Annalen der Physik;
Zeitschrift fur Physik

\subsection*{Consulting Activities}
\noindent
\begin{tabular}{lll}
    1939-1945&
  Scientific Advisor&
  United States Government (Manhattan Project advisory)\\
    \end{tabular}
  

%%End Service--------------------------------

\section*{Public Health Practice and Communication}

(illustrative selection)

\subsection*{Presentations to policy-makers and other stakeholders}

\setlength{\extrarowheight}{.75em}
\begin{longtable}[l]{lp{5in}}   
November 1945&
\parbox[t]{5in}{
  Atomic energy and its implications for international policy%
  \textit{, United States Senate Special Committee on Atomic Energy%
    }%
    : %
    Testified on the need for international control of atomic energy%
    .%
  }\\
October 1947&
\parbox[t]{5in}{
  The case for a world government%
  \textit{, United Nations General Assembly Committee%
    }%
    : %
    Presented argument for supranational governance to prevent nuclear
war%
    .%
  }\\
\end{longtable}
\setlength{\extrarowheight}{0em}

\subsection*{Consultations with policy-makers and other
  stakeholders}

\setlength{\extrarowheight}{.75em}
\begin{longtable}[l]{lp{5in}}   
1941-1945&
\parbox[t]{5in}{
    \textit{Manhattan Project Advisory Consultation%
    , United States War Department%
    }%
    : %
    Provided advisory consultation on theoretical physics related to
nuclear fission%
    .%
  }\\
1946-1950&
\parbox[t]{5in}{
    \textit{Atomic Energy Commission Advisory%
    , United States Atomic Energy Commission%
    }%
    : %
    Advised on peaceful applications of atomic energy and international
arms control%
    .%
  }\\
\end{longtable}
\setlength{\extrarowheight}{0em}

\subsection*{Research finding dissemination through media appearances
  and other communication venues}


\setlength{\extrarowheight}{.75em}
\begin{longtable}[l]{lp{5in}}   
1919&
\parbox[t]{5in}{
    General relativity confirmed by solar eclipse observations: %
    The Times (London); The New York Times.%    
  }\\
1939&
\parbox[t]{5in}{
    Letter to President Roosevelt regarding atomic energy: %
    The New York Times.%    
  }\\
\end{longtable}
\setlength{\extrarowheight}{0em}


\subsection*{Other practice activities}

\setlength{\extrarowheight}{.75em}
\begin{longtable}[l]{lp{5in}}   
1946-1948&
\parbox[t]{5in}{
    \textit{Chairman, Emergency Committee of Atomic Scientists%
    , Emergency Committee of Atomic Scientists%
    }%
    : %
    Led fundraising and public education campaign on dangers of nuclear
weapons and need for international control%
    .%
  }\\
1955&
\parbox[t]{5in}{
    \textit{Russell-Einstein Manifesto%
    , International scientific community%
    }%
    : %
    Co-authored manifesto with Bertrand Russell calling for peaceful
resolution of international conflicts and warning of nuclear war
dangers%
    .%
  }\\
\end{longtable}
\setlength{\extrarowheight}{0em}

%%%----------------END PRACTICE ACTIVITIES

\end{document}